\section{Case Study: Multi-Modal Query on the Karpathy Corpus}
\label{sec:case}

To illustrate the practical value of multi-modal graph traversal, we trace a
concrete query through Graphify on the 52-file karpathy corpus (nanoGPT,
minGPT, micrograd + 5 research papers + 4 images).

\subsection{Query: ``how does the attention mechanism work''}

\textbf{Baseline (grep / file reading).} A Bash grep for \texttt{attention} in
the corpus returns matches across 9 files spanning 14,700 tokens --- over 10$\times$
the context needed and with no structural ordering. The agent must read each file
sequentially to reconstruct the relationship between the code implementation and
the paper describing it.

\textbf{Graphify BFS traversal.} The hub-aware BFS seeds on three nodes scored
highest by Eq.~\eqref{eq:score}: \texttt{CausalSelfAttention} (minGPT code),
\texttt{FlashAttention: Fast and Memory-Efficient Exact Attention} (PDF), and
\texttt{Online Softmax for Attention} (PDF). Expanding to depth 2 with $\tau =
p_{99}(\deg)$ produces the subgraph in Table~\ref{tab:casestudy}: 2,836 tokens
covering code, paper, and cross-modal edges that the grep baseline cannot express.

\begin{table}[h]
\centering
\caption{Subgraph returned for ``how does the attention mechanism work.''
         Cross-modal edges (\textit{italics}) span code and research papers.}
\label{tab:casestudy}
\begin{tabular}{lll}
\toprule
\textbf{Node} & \textbf{Source} & \textbf{Relation} \\
\midrule
CausalSelfAttention      & minGPT/model.py      & seed     \\
\quad .forward()         & minGPT/model.py      & contains \\
\quad CfgNode            & minGPT/model.py      & uses     \\
\quad \textit{Causal Self-Attention}   & \textit{minGPT/model.py}
                         & \textit{concept.related} \\
\quad\quad \textit{Attention Is All You Need}  & \textit{PDF (Vaswani 2017)}
                         & \textit{cites} \\
FlashAttention (PDF)     & flashattention.pdf   & seed     \\
\quad Tiling Technique   & flashattention.pdf   & implements \\
\quad Block-Sparse Attn  & flashattention.pdf   & implements \\
\quad \textit{CausalSelfAttention}     & \textit{nanoGPT/model.py}
                         & \textit{implements} \\
Online Softmax           & flashattention.pdf   & seed     \\
\quad Milakov \& Gimelshein & flashattention.pdf & cites    \\
\bottomrule
\end{tabular}
\end{table}

The graph surfaces a cross-modal edge that neither grep nor static analysis
can produce: \texttt{CausalSelfAttention} in nanoGPT's code is connected via
an INFERRED \texttt{implements} edge to the FlashAttention algorithm in the PDF,
which was not imported or called --- it was inferred by the LLM semantic pass
from docstrings referencing the paper. This is the structural knowledge that
justifies the graph abstraction: once built, it answers not just ``where is
attention implemented?'' but ``which paper describes the algorithm this code
implements and what did that paper cite?'' --- across file modalities, in 2,836
tokens instead of 14,700.

\subsection{Surprising Cross-Repository Connections}

A second class of graph-exclusive insight is cross-repository relationships that
static analysis misses because imports do not cross repo boundaries. On the
karpathy corpus, Graphify's semantic pass infers:

\begin{itemize}
  \item \texttt{from\_pretrained()} in nanoGPT \texttt{--calls-->}
        \texttt{get\_default\_config()} in minGPT (\textsc{Inferred}, c=0.87):
        nanoGPT's weight loading follows the same interface contract as minGPT,
        a cross-repo architectural dependency invisible to any single-repo call
        graph.
  \item \texttt{CausalSelfAttention} in minGPT \texttt{--conceptually\_related\_to-->}
        \texttt{CausalSelfAttention} in nanoGPT (\textsc{Inferred}, c=0.93):
        Andrej Karpathy reimplemented the same module; the graph names both and
        links them, making the lineage queryable.
  \item \texttt{GPT Language Model (minGPT)} \texttt{--conceptually\_related\_to-->}
        \texttt{GPT Model Class} in nanoGPT (\textsc{Inferred}): identifies the
        successor/predecessor relationship between the two repos.
\end{itemize}

\noindent These edges are tagged \textsc{Inferred} with high confidence
($c \geq 0.87$), and surfaced in the \texttt{GRAPH\_REPORT.md} ``Surprising
Connections'' section, giving the developer a structured starting point for
cross-repo architecture review without any manual curation.

\section{Discussion}
\label{sec:discussion}

\textbf{When graph traversal beats file reading.} Token reduction scales with
corpus size: on 6-file corpora the graph provides structural clarity but no
compression. At 52 files, the BFS subgraph stays roughly constant (1,726 mean
tokens) while naive stuffing scales linearly, yielding 71.5$\times$ savings. The
practical threshold in our experiments is roughly 20--30 files: below that,
graph traversal is an organisational tool; above it, the token economics become
compelling.

\textbf{Confidence tagging enables trust calibration.} A developer querying the
graph can see whether an edge was syntactically verified (\textsc{Extracted}) or
LLM-inferred (\textsc{Inferred}). In our corpora, cross-repo and cross-modal
edges are almost exclusively \textsc{Inferred}; within-file structural edges are
almost exclusively \textsc{Extracted}. This separation gives the developer a
reliable signal: EXTRACTED edges can be trusted unconditionally; INFERRED edges
are high-confidence hypotheses worth verifying.

\textbf{Threats to validity.} (1) \emph{Token reduction} is measured against a
naive full-corpus baseline; a fairer comparison against an optimised agentic
reader (which only reads files it selects) would show smaller reduction ratios,
as the agentic baseline already targets relevant files. (2) \emph{Confidence
distribution} is system-asserted, not validated against a human annotation; an
EXTRACTED edge can still be incorrect if the AST walker resolves an import to the
wrong symbol. (3) \emph{Corpora} are Python-dominant; token reduction on
multi-language codebases may differ due to varying file sizes and import
graph density. (4) \emph{Parallel speedup} is modest on our test machine (1.38$\times$
on 160 files); larger codebases with more CPU-intensive per-file parsing (e.g.,
C++ with complex templates) are expected to show higher speedups.


\subsection{Community Structure as Subsystem Map}

The 53 communities detected on the karpathy corpus by topology-only Leiden
partitioning correspond to recognisable software subsystems without any
domain-specific tuning: ``nanoGPT Model Architecture'' (12 nodes), ``FlashAttention
Paper'' (21 nodes), ``BPE Tokenizer'' (8 nodes), and ``micrograd Autograd
Engine'' (6 nodes), among others. Community cohesion scores range from 0.10
to 0.21, with isolate communities (single-node) at 1.0 by definition.

Critically, the FlashAttention paper (5 PDFs, 4 images) forms its own community
entirely separate from the nanoGPT code community, yet a cross-community edge
\texttt{CausalSelfAttention} $\rightarrow$ \texttt{FlashAttention Algorithm}
spans them. This inter-community edge is flagged in the ``Surprising Connections''
report, directing the developer's attention to the code-paper dependency that
is otherwise invisible without reading both sources. The graph's community
structure thus acts as a high-level subsystem map, while cross-community edges
flag non-obvious architectural dependencies.
